
\documentclass{article}
\usepackage{colortbl}
\usepackage{makecell}
\usepackage{multirow}
\usepackage{supertabular}

\begin{document}

\newcounter{utterance}

\centering \large Interaction Transcript for game `eqbench', experiment `eq\_bench', episode 131 with reflexion{-}Qwen3.5{-}27B.
\vspace{24pt}

{ \footnotesize  \setcounter{utterance}{1}
\setlength{\tabcolsep}{0pt}
\begin{supertabular}{c@{$\;$}|p{.15\linewidth}@{}p{.15\linewidth}p{.15\linewidth}p{.15\linewidth}p{.15\linewidth}p{.15\linewidth}}
    \# & \multicolumn{2}{c}{Player} && \multicolumn{2}{c}{Game Master} \\
    \hline

    \theutterance \stepcounter{utterance}  
    & & & \multicolumn{4}{p{0.6\linewidth}}{
        \cellcolor[rgb]{0.9,0.9,0.9}{
            \makecell[{{p{\linewidth}}}]{
                \texttt{\tiny{[P1$\langle$GM]}}
                \texttt{You have attempted this same task before and did not succeed. These are the plans you wrote after those attempts.} \\
\texttt{Trial \#1: New plan: My previous attempt failed because I underestimated the intensity of Sam's skepticism and overestimated the likelihood of him feeling vindicated. Sam's dialogue ("And I'm supposed to just accept it because you've decided to say sorry?") explicitly rejects the apology, indicating that "Pacified" should be near zero and "Skeptical" should be the dominant emotion (likely 10). "Vindicated" is weak because Sam is still angry, not triumphant. Next time, I will assign a score of 10 to Skeptical, 0 to Pacified and Joyful, and a low score (e.g., 1 or 2) to Vindicated, ensuring the scores accurately reflect the character's explicit rejection of the apology.} \\
\\ 
\texttt{Trial \#2: New plan: My previous attempt failed because I assigned a non{-}zero score (2) to "Vindicated," which misinterpreted Sam's emotional state. Sam is not feeling triumphant or justified in a victory; rather, he is actively rejecting Arthur's apology and expressing deep distrust. The dialogue "And I'm supposed to just accept it because you've decided to say sorry?" confirms that Sam feels nothing but doubt. To correct this, I will assign a score of 10 to "Skeptical" as the dominant emotion, and strictly assign 0 to "Pacified," "Vindicated," and "Joyful," ensuring the scores reflect a complete lack of resolution or triumph in Sam's response.} \\
\\ 
\texttt{Trial \#3: New plan: My previous attempts failed because I incorrectly assigned a non{-}zero score to "Vindicated" in Trial 1 and then over{-}corrected by rigidly assigning 0 to everything else in Trial 2, potentially missing the nuance that Sam feels a degree of vindication in his anger (e.g., "Just like your speech at my wedding" implies he feels justified in his past grievances). The core error was not fully analyzing the specific text "Just like your speech at my wedding," which suggests Sam feels his past negative experiences were validated by Arthur's current behavior. To fix this, I will assign 10 to "Skeptical" as the primary emotion due to the explicit rejection of the apology. I will also assign a moderate score (e.g., 5) to "Vindicated" to reflect Sam's feeling of being right about Arthur's character flaws, while keeping "Pacified" and "Joyful" at 0. This balances the dominant skepticism with the secondary feeling of justified grievance.} \\
\\ 
\\ 
\\ 
\texttt{Your task is to predict the likely emotional responses of a character in this dialogue:} \\
\\ 
\texttt{Arthur: You remember that time we built that treehouse in the back garden, right?} \\
\texttt{Sam: What's that got to do with anything?} \\
\texttt{Arthur: I'm just saying, we were a team, you and me. We built something out of nothing. } \\
\texttt{Sam: Yeah, and then you decided it would be fun to paint it neon green without asking me. } \\
\texttt{Arthur: It was a creative decision, Sammy. } \\
\texttt{Sam: It was a stupid decision, Arthur. Just like your speech at my wedding.} \\
\texttt{Arthur: That's not fair, Sam. I messed up, I know. But I'm trying to apologize here.} \\
\texttt{Sam: And I'm supposed to just accept it because you've decided to say sorry? } \\
\texttt{Arthur: No, not because I've decided to. Because I mean it, Sam. I'm sorry.} \\
\texttt{{[}End dialogue{]}} \\
\\ 
\texttt{At the end of this dialogue, Sam would feel...} \\
\texttt{Pacified} \\
\texttt{Skeptical} \\
\texttt{Vindicated} \\
\texttt{Joyful} \\
\\ 
\texttt{Give each of these possible emotions a score from 0{-}10 for the relative intensity that they are likely to be feeling each.} \\
\\ 
\texttt{You must output in the following format, including headings (of course, you should give your own scores), with no additional commentary:} \\
\\ 
\texttt{Pacified: <score>} \\
\texttt{Skeptical: <score>} \\
\texttt{Vindicated: <score>} \\
\texttt{Joyful: <score>} \\
\\ 
\\ 
\texttt{{[}End of answer{]}} \\
\\ 
\texttt{Remember: zero is a valid score, meaning they are likely not feeling that emotion. You must score at least one emotion > 0.} \\
\\ 
\texttt{Your answer:} \\
            }
        }
    }
     \\ \\

    \theutterance \stepcounter{utterance}  
    & \multicolumn{4}{p{0.6\linewidth}}{
        \cellcolor[rgb]{0.561,0.737,0.561}{
            \makecell[{{p{\linewidth}}}]{
                \texttt{\tiny{[P1$\rangle$GM]}}
                \texttt{Pacified: 0} \\
\texttt{Skeptical: 10} \\
\texttt{Vindicated: 6} \\
\texttt{Joyful: 0} \\
            }
        }
    }
    & & \\ \\

    \theutterance \stepcounter{utterance}  
    & & & \multicolumn{2}{p{0.3\linewidth}}{
        \cellcolor[rgb]{0.9,0.9,0.9}{
            \makecell[{{p{\linewidth}}}]{
                \texttt{\tiny{[GM$|$GM]}}
                \texttt{Pacified: 0} \\
\texttt{Skeptical: 10} \\
\texttt{Vindicated: 6} \\
\texttt{Joyful: 0} \\
            }
        }
    }
    & & \\ \\

    \theutterance \stepcounter{utterance}  
    & & & \multicolumn{2}{p{0.3\linewidth}}{
        \cellcolor[rgb]{0.9,0.9,0.9}{
            \makecell[{{p{\linewidth}}}]{
                \texttt{\tiny{[GM$|$GM]}}
                \texttt{Pacified: 4} \\
\texttt{Skeptical: 7} \\
\texttt{Vindicated: 3} \\
\texttt{Joyful: 0} \\
            }
        }
    }
    & & \\ \\

    \theutterance \stepcounter{utterance}  
    & & & \multicolumn{2}{p{0.3\linewidth}}{
        \cellcolor[rgb]{0.9,0.9,0.9}{
            \makecell[{{p{\linewidth}}}]{
                \texttt{\tiny{[GM$|$GM]}}
                \texttt{game\_result = LOSE} \\
            }
        }
    }
    & & \\ \\

    \theutterance \stepcounter{utterance}  
    & & \multicolumn{4}{p{0.6\linewidth}}{
        \cellcolor[rgb]{0.95,0.95,0.95}{
            \makecell[{{p{\linewidth}}}]{
                \texttt{\tiny{[P1$\leftrightarrow$P]}}
                \texttt{Post{-}episode Reflexion (not sent to GM):} \\
\texttt{New plan: My previous attempts failed because I oscillated between underestimating and overestimating the "Vindicated" score without a consistent analytical anchor. The core error was not strictly differentiating between "feeling right about a past grievance" (Vindicated) and "distrusting a current apology" (Skeptical). Sam's line "Just like your speech at my wedding" confirms he feels his past judgment was correct, justifying a non{-}zero "Vindicated" score, but his rhetorical question "And I'm supposed to just accept it...?" confirms he is not "Pacified" or "Joyful." To succeed, I will assign 10 to "Skeptical" as the dominant reaction to the failed apology. I will assign a moderate score (e.g., 5) to "Vindicated" to acknowledge his feeling of being proven right about Arthur's character, while strictly keeping "Pacified" and "Joyful" at 0. This ensures the output reflects both the active rejection of the apology and the validation of his past anger.} \\
            }
        }
    }
    & \\ \\

\end{supertabular}
}

\end{document}
